% Generated draft: requires mathematical and source review.
\documentclass[11pt]{article}
\usepackage[margin=25mm]{geometry}
\usepackage{fontspec}
\setmainfont{Latin Modern Roman}
\usepackage{amsmath,amssymb,mathtools,unicode-math}
\setmathfont{Latin Modern Math}
\usepackage{enumitem,needspace,xurl,hyperref}
\setlength{\parindent}{0pt}
\setlength{\parskip}{5pt}
\setlength{\emergencystretch}{3em}
\newcommand{\N}{\mathbb{N}}
\newcommand{\Z}{\mathbb{Z}}
\newcommand{\Q}{\mathbb{Q}}
\newcommand{\R}{\mathbb{R}}
\newcommand{\C}{\mathbb{C}}
\newcommand{\F}{\mathbb{F}}
\newcommand{\A}{\mathbb{A}}
\newcommand{\PP}{\mathbb{P}}
\newcommand{\E}{\mathbb{E}}
\newcommand{\eps}{\varepsilon}
\newcommand{\dd}{\,\mathrm{d}}
\begin{document}
\section*{UnsolvedMath ID: 30005696}
\textbf{Source problem code:} catalogue entry
\subsection{Definitions and mathematical statement}
Let \(G\) be a totally disconnected locally compact group and let \(O\leq G\) be a compact open subgroup. Let \(\mu\) be a left Haar measure on \(G\), and write \(\mu_O=\mu/\mu(O)\) for the normalized Haar measure satisfying \(\mu_O(O)=1\). Choose a set \(R\subseteq G\) of representatives of the double cosets \(O\backslash G/O\). For each integer \(n\geq 1\), define
\[
 a_n(G,O)=\#\{r\in R:\mu(OrO)=n\mu(O)\}.
\]
The group has \emph{bounded coset growth with respect to \(O\)} in the terminology of the supplied excerpt if \(a_n(G,O)<\infty\) for every \(n\geq1\). This condition does not assert a uniform bound on the coefficients. When it holds, define the formal Dirichlet series
\[
 \zeta_{G,O}(s)=\sum_{n=1}^{\infty}a_n(G,O)n^{-s}.
\]
The coefficient definition is independent of the chosen representatives and of scaling \(\mu\). The supplied excerpt states that bounded coset growth with respect to one compact open subgroup implies bounded coset growth with respect to every compact open subgroup.

The hypotheses in the cleaned record are that \(G\) is unimodular, rationally of type \(\mathrm{FP}_{\infty}\), and has finite rational cohomological dimension. Under the usual discrete-module convention for totally disconnected locally compact groups, rational type \(\mathrm{FP}_{\infty}\) means that the trivial module \(\Q\) has a projective resolution whose terms are finitely generated in every degree, in the category of discrete \(\Q[G]\)-modules with continuous \(G\)-action. Finite rational cohomological dimension means that \(\Q\) has finite projective dimension in this category. The supplied excerpt does not specify this category, so this convention requires confirmation against the report.

\emph{Requested characterization.} Determine precisely which pairs \((G,O)\) in the intended source class, subject to these hypotheses, satisfy
\[
 \chi_G=\widehat\zeta_{G,O}(-1)\,\mu_O.
\]
Here \(\chi_G\) denotes the Euler characteristic intended by the source, apparently as an element of a space of Haar measures, and \(\widehat\zeta_{G,O}\) denotes the analytic object intended by the cleaned record. If the hat denotes meromorphic continuation of the displayed Dirichlet series, the equality requires a well-defined finite value at \(s=-1\); the record supplies no regularization prescription for a pole there.

\emph{Unresolved formulation qualifications.} The original excerpt calls the pairs \emph{merohoric}, whereas the cleaned statement calls them \emph{meromorphic}. Neither term is defined in the supplied material, and their equivalence cannot be assumed. The excerpt also omits the original conditions labelled (i)--(iii), the definition of \(\chi_G\), and the definition of the hatted zeta function. Consequently, the displayed characterization preserves the supplied target but remains conditional on recovering these definitions and confirming the normalization of \(\mu_O\).
\subsection{Short English statement}
For a totally disconnected locally compact group \(G\) and a compact open subgroup \(O\), count double cosets by their Haar measure relative to \(O\) and form the resulting Dirichlet series. The problem asks which intended merohoric or meromorphic pairs, with \(G\) unimodular, rationally of type \(\mathrm{FP}_{\infty}\), and of finite rational cohomological dimension, recover the Euler characteristic by evaluating the intended continuation at \(-1\). The source does not supply enough definitions to identify the pair class or the analytic evaluation unambiguously.
\subsection{Sources}
\begin{itemize}
\item Local UnsolvedMath registry, permanent ID \(30005696\), exact source problem code \texttt{catalogue entry}, title: \emph{Euler Characteristics from Formal Dirichlet Series of Merohoric Pairs}. The supplied record belongs to the Oberwolfach Reports Open Problems collection. No individual proposer is identified.
\item Supplied report citation: \emph{Mini-Workshop: Homological aspects of TDLC-groups (2024)}, Oberwolfach Reports, pp.\ 3103--3138, DOI \texttt{10.4171/owr/2023/55}. Supplied link: \url{https://doi.org/10.4171/owr/2023/55}. The registry separately records proposed year 2023; the citation's year discrepancy is unresolved here.
\item Supplied literature source, labelled ``Oberwolfach report'': \url{https://ems.press/content/serial-article-files/48166}.
\item Literature assessment dated exactly \texttt{2026-08-22}: the supplied curation labels the problem open, reports proofs for several Bruhat-decomposition families, reports counterexamples among general meromorphic pairs, and states that no characterization of all pairs is known. This is supplied dated curation, not an independently verified literature assessment.
\item The supplied excerpt reports an affirmative answer to ``Question 6'' for \((\widehat X_D(F),\mathrm{Iw})\) for every crystallographic Coxeter group \((W,S)\), referring to its Example 1(E), and for \(G=X(E)\), with \(X\) a Chevalley group scheme, \(E\) a totally disconnected locally compact field, and \(O\) parahoric, referring to its reference [5]. Those internal references are not supplied in full. It also attributes to B.\ Marchionna examples of merohoric pairs failing its equation (4), and many satisfying (4) for groups without Bruhat decomposition. Equation (4) is absent, so its relationship to the requested identity is unverified.
\item The registry labels its statement \texttt{corrected\_verified} and says it was checked against the official report. These are supplied provenance claims; no independent source consultation or mathematical verification is asserted in this draft.
\end{itemize}
\end{document}
